For the case $\lij = \dijEq$ the second term of (\ref{eq:kmatFrustrated}) vanishes and the first term gives the familiar result. The contribution of the second term is due to frustration. Since $\kmat$ depends on frustration, protein dynamics for the frustrated case should be different than for the non-frustrated case. For instance, normal modes, which are the eigenvectors of $\kmat$ will depend on the departure of $\lij/\dijEq$ from $1$. 



\subsection{?}

At the potential energy minimum we have $\rvectorEq$ and $\dijEq = ||\rjEq - \riEq||$. 

All ENM developed so far have assumed $\lij = \dijEq$ in (\ref{eq:vENM}). This means that such modes assume that at equilibrium, all springs are completely relaxed. This choice guarantees that the minimum of the potential is at the native conformation of the protein considered. However, there is frustration in proteins, and such frustration could be modeled by choosing $\lij$ with the weaker requirement that the minimum is at $\rvectorEq$ without necessarily having all springs relaxed. In general 
\begin{equation} \label{eq:VminFrustrated}
V(\rvectorEq) = V^0 + \frac{1}{2} \sum_{i<j} \kij (\dijEq - \lij)^2
\end{equation}
From this it is clear that for the non-frustrated case the minimum is $V^0$ and that the second term represents the frustration.